%Daten zur Institution

%\input{Dozentdaten}

%\renewcommand{\fachbereich}{Fachbereich}

%\renewcommand{\dozent}{Prof. Dr. . }

%Klausurdaten

\renewcommand{\klausurgebiet}{ }

\renewcommand{\klausurtyp}{ }

\renewcommand{\klausurdatum}{ . 20}

\klausurvorspann {\fachbereich} {\klausurdatum} {\dozent} {\klausurgebiet} {\klausurtyp}

%Daten für folgende Punktetabelle

\renewcommand{\aeins}{  3  }

\renewcommand{\azwei}{  3   }

\renewcommand{\adrei}{  0  }

\renewcommand{\avier}{  0  }

\renewcommand{\afuenf}{  6  }

\renewcommand{\asechs}{  6  }

\renewcommand{\asieben}{  0  }

\renewcommand{\aacht}{  6  }

\renewcommand{\aneun}{  11  }

\renewcommand{\azehn}{  5  }

\renewcommand{\aelf}{  5  }

\renewcommand{\azwoelf}{  3  }

\renewcommand{\adreizehn}{  7   }

\renewcommand{\avierzehn}{  2  }

\renewcommand{\afuenfzehn}{  0  }

\renewcommand{\asechzehn}{  6  }

\renewcommand{\asiebzehn}{  63  }

\renewcommand{\aachtzehn}{    }

\renewcommand{\aneunzehn}{    }

\renewcommand{\azwanzig}{    }

\renewcommand{\aeinundzwanzig}{    }

\renewcommand{\azweiundzwanzig}{    }

\renewcommand{\adreiundzwanzig}{    }

\renewcommand{\avierundzwanzig}{    }

\renewcommand{\afuenfundzwanzig}{    }

\renewcommand{\asechsundzwanzig}{    }

\punktetabellesechzehn

\klausurnote

\newpage

\setcounter{section}{K}

\inputaufgabegibtloesung
{3}
{

Definiere die folgenden
\zusatzklammer {kursiv gedruckten} {} {} Begriffe.
\aufzaehlungsechs{Eine
\stichwort {geodätische Kurve} {}
auf einer differenzierbaren Hyperfläche

\mavergleichskette
{\vergleichskette
{ Y
}
{ \subseteq }{ \R^n
}
{  }{
}
{    }{
}
{   }{
}
}
{}{}{.}

}{Die \stichwort {Rotationsmenge} {}
\zusatzklammer {um die $x$-Achse} {} {}
zu einer Teilmenge

\mavergleichskette
{\vergleichskette
{ T
}
{ \subseteq }{  \R \times \R_{\geq 0}
}
{  }{
}
{    }{
}
{   }{
}
}
{}{}{.}

}{Die
\stichwort {Tangentialabbildung} {}
\maabbdisp {T (\varphi)} {TM} {TN
} {}
zu einer
\definitionsverweis {differenzierbaren Abbildung}{}{}
\maabbdisp {\varphi} { M } { N
} {}
zwischen
\definitionsverweis {differenzierbaren Mannigfaltigkeiten}{}{}
\mathkor {} {M} {und} {N} {.}

}{Ein
\stichwort {orientierungstreuer} {}
Kartenwechsel auf einer
\definitionsverweis {differenzierbaren Mannigfaltigkeit}{}{}
$M$.

}{Die \stichwort {zurückgezogene Differentialform} {} $\varphi^* \omega$ zu einer Differentialform
\mathl{\omega \in
{ \mathcal E }^{ k } (  M   )}{} bezüglich einer stetig differenzierbaren Abbildung
\maabb {\varphi} { L } { M
} {}
zwischen zwei differenzierbaren Mannigfaltigkeiten
\mathkor {} {L} {und} {M} {.}

}{Ein
\stichwort {metrischer} {}
linearer Zusammenhang auf dem Tangentialbündel einer riemannschen Mannigfaltigkeit $M$.
}

}
{} {}

\inputaufgabegibtloesung
{3}
{

Formuliere die folgenden Sätze.
\aufzaehlungdrei{Der
\stichwort {Satz über die Lösung von gewöhnlichen Differentialgleichungen auf einer differenzierbaren Hyperfläche} {}

\mavergleichskette
{\vergleichskette
{  Y
}
{ \subseteq }{  \R^n
}
{  }{
}
{    }{
}
{   }{
}
}
{}{}{.}}{Die Formel für die Berechnung des kanonischen Volumens auf einer riemannschen Mannigfaltigkeit in einer Karte.}{Der Satz über die Partition der Eins.}

}
{} {}

\inputaufgabe
{0}
{

}
{} {}

\inputaufgabe
{0}
{

}
{} {}

\inputaufgabegibtloesung
{6}
{

Bestimme für die durch

\mavergleichskettedisp
{\vergleichskette
{ x^2+y^2+2z^2
}
{ =} { 4
}
{ } {
}
{ } {
}
{ } {
}
}
{}{}{}
gegebene Fläche $Y$ und den Punkt

\mavergleichskettedisp
{\vergleichskette
{ P
}
{ =} { \left(  1  , \,   1  , \,   1                 \right)
}
{ \in} { Y
}
{ } {
}
{ } {
}
}
{}{}{}
eine Diagonalmatrix für die
\definitionsverweis {Weingartenabbildung}{}{}
$L_P$.

}
{} {}

\inputaufgabegibtloesung
{6}
{

Beweise den Isometriesatz für den Paralleltransport auf einer Hyperfläche

\mavergleichskette
{\vergleichskette
{ Y
}
{ \subseteq }{ \R^n
}
{  }{
}
{    }{
}
{   }{
}
}
{}{}{.}

}
{} {}

\inputaufgabe
{0}
{

}
{} {}

\inputaufgabegibtloesung
{6}
{

Es sei $M$ eine kompakte topologische $d$-dimensionale Mannigfaltigkeit,

\mavergleichskette
{\vergleichskette
{ d
}
{ \geq }{ 1
}
{  }{
}
{    }{
}
{   }{
}
}
{}{}{.}
Zeige, dass es eine beschränkte offene Teilmenge

\mavergleichskette
{\vergleichskette
{ U
}
{ \subseteq }{ \R^d
}
{  }{
}
{    }{
}
{   }{
}
}
{}{}{}
und eine stetige surjektive Abbildung
\maabbdisp {\varphi} { U } { M
} {}
gibt.

}
{} {}

\inputaufgabegibtloesung
{11 (1+3+1+2+4)}
{

Wir betrachten die Menge

\mavergleichskettedisp
{\vergleichskette
{N
}
{ =} { \left\{ A= \begin{pmatrix}  a   &   b   \\  c  & d \end{pmatrix}   \mid   A \text{ ist nilpotent}         \right\}
}
{ \subseteq} {\R^4
}
{ } {
}
{ } {
}
}
{}{}{}
der reellen nilpotenten
$2\times 2$-\definitionsverweis {Matrizen}{}{}
sowie die Menge

\mavergleichskettedisp
{\vergleichskette
{M
}
{ =} { N \setminus \{ \begin{pmatrix}  0   &   0  \\  0  &  0 \end{pmatrix}  \}
}
{ } {
}
{ } {
}
{ } {
}
}
{}{}{.}

a) Ist $N$ zusammenhängend?

b) Zeige, dass $M$ eine
\definitionsverweis {abgeschlossene Untermannigfaltigkeit}{}{}
einer offenen Teilmenge

\mavergleichskette
{\vergleichskette
{ G
}
{ \subseteq }{ \R^4
}
{  }{
}
{    }{
}
{   }{
}
}
{}{}{}
ist.

c) Bestimme die
\definitionsverweis {Dimension}{}{}
von $M$.

d) Ist $M$ zusammenhängend?

e) Überdecke
\mathl{M}{} mit expliziten topologischen Karten.

}
{} {}

\inputaufgabe
{5}
{

Sagen Sie etwas Schlaues zum Thema Möbiusband.

}
{} {}

\inputaufgabegibtloesung
{5}
{

Wir betrachten die Abbildung
\maabbeledisp {\varphi} { \R^3 } { \R^2
} { (x,y,z) } { \left(  xyz^2,xy-z^3  \right)
} {.}
Berechne die Matrix der Abbildung
\maabbdisp {\bigwedge^2 T_P (\varphi)} { \bigwedge^2 T_P \R^3 } { \bigwedge^2 T_{\varphi(P)}  \R^2
} {}
im Punkt

\mavergleichskette
{\vergleichskette
{ P
}
{ = }{ (1,3,5)
}
{  }{
}
{    }{
}
{   }{
}
}
{}{}{}
bezüglich einer geeigneten Basis.

}
{} {}

\inputaufgabegibtloesung
{3}
{

Es seien

\mavergleichskette
{\vergleichskette
{ W
}
{ \subseteq }{ \R^m
}
{  }{
}
{    }{
}
{   }{
}
}
{}{}{}
und

\mavergleichskette
{\vergleichskette
{ U
}
{ \subseteq }{ \R^n
}
{  }{
}
{    }{
}
{   }{
}
}
{}{}{}
\definitionsverweis {offene Teilmengen}{}{}
und sei
\maabbdisp {\psi} { W } { U
} {}
eine
\definitionsverweis {stetig differenzierbare}{}{}
\definitionsverweis {Abbildung}{}{.}
Es sei
\maabbdisp {f} { U } {\R
} {}
eine stetig differenzierbare Funktion. Folgere aus der
[[Totale Differenzierbarkeit/K/Kettenregel/Fakt|Kettenregel]],
dass

\mavergleichskettedisp
{\vergleichskette
{ d (\psi^*f)
}
{ =} { \psi^*(df)
}
{ } {
}
{ } {
}
{ } {
}
}
{}{}{}
gilt, wobei $\psi^*$ das
\definitionsverweis {Zurückziehen von Differentialformen}{}{}
bezeichnet.

}
{} {}

\inputaufgabegibtloesung
{7}
{

Es sei $M$ eine
\definitionsverweis {differenzierbare Mannigfaltigkeit}{}{}
mit
\definitionsverweis {abzählbarer Basis der Topologie}{}{}
und sei $\omega$ eine
\definitionsverweis {positive Volumenform}{}{}
auf $M$. Es sei
\maabbdisp {\varphi} { L } { M
} {}
ein
\definitionsverweis {Diffeomorphismus}{}{}
mit der Mannigfaltigkeit $L$ und

\mavergleichskette
{\vergleichskette
{ S
}
{ \subseteq }{ L
}
{  }{
}
{    }{
}
{   }{
}
}
{}{}{}
eine
\definitionsverweis {messbare Teilmenge}{}{.}
Zeige

\mavergleichskettedisp
{\vergleichskette
{ \int_S \varphi^* \omega
}
{ =} { \int_{ \varphi(S)} \omega
}
{ } {
}
{ } {
}
{ } {
}
}
{}{}{.}

}
{} {}

\inputaufgabegibtloesung
{2}
{

Wir betrachten auf dem
\definitionsverweis {trivialen Vektorbündel}{}{}
\maabbdisp {p} {\R^2 \times \R} { \R^2
} {}
vom
\definitionsverweis {Rang}{}{}
$1$ über $\R^2$ den
\definitionsverweis {linearen Zusammenhang}{}{,}
der durch die
\definitionsverweis {Christoffelsymbole}{}{}

\mavergleichskette
{\vergleichskette
{ \Gamma_1 (x,y)
}
{ = }{ x^5-x^2y^2+3y^4
}
{  }{
}
{    }{
}
{   }{
}
}
{}{}{}
und

\mavergleichskette
{\vergleichskette
{ \Gamma_2 (x,y)
}
{ = }{ 7x^3 +xy^2-4y^5
}
{  }{
}
{    }{
}
{   }{
}
}
{}{}{}
gegeben sei. Berechne den
\definitionsverweis {Krümmungsoperator}{}{}
$R \left(  \partial_1, \partial_2  \right)$.

}
{} {}

\inputaufgabe
{0}
{

}
{} {}

\inputaufgabegibtloesung
{6}
{

Beweise den Satz über die Retraktionen auf den Rand auf Mannigfaltigkeiten.

}
{} {}

 [[Kategorie:Latexseite]]
